\documentclass[11pt,letterpaper]{article}
\usepackage[T1]{fontenc}
\usepackage{lmodern,amsmath,amssymb,amsthm,mathtools,booktabs,longtable,array}
\usepackage[margin=0.9in,headheight=14pt]{geometry}
\usepackage{microtype,xcolor,fancyhdr,listings,enumitem}
\usepackage[colorlinks=true,linkcolor=teal!55!black,urlcolor=teal!55!black,citecolor=teal!55!black]{hyperref}
\hypersetup{pdftitle={iPRA: smaller checkers and the independence of progress},pdfauthor={Prepared for Vladimir Reshetnikov},pdfsubject={Intuitionistic primitive recursive arithmetic; explicit proof presentation; coherent ultrapower countermodel}}
\newtheorem{theorem}{Theorem}[section]
\newtheorem{lemma}[theorem]{Lemma}
\newtheorem{proposition}[theorem]{Proposition}
\newtheorem{corollary}[theorem]{Corollary}
\theoremstyle{definition}
\newtheorem{definition}[theorem]{Definition}
\newtheorem{remark}[theorem]{Remark}
\newcommand{\ipra}{\mathsf{iPRA}}
\newcommand{\I}{\mathsf I}
\newcommand{\Q}{\mathsf Q}
\newcommand{\D}{\mathsf D}
\newcommand{\K}{\mathcal K}
\newcommand{\Km}{\mathcal K^{-}}
\newcommand{\G}{\Gamma}
\newcommand{\llbracket}{[\![}
\newcommand{\rrbracket}{]\!]}
\newcommand{\file}[1]{{\small\nolinkurl{#1}}}
\newcommand{\srcbase}{https://github.com/VladimirReshetnikov/ProveIt/blob/6ffc49f55d355b548f7a8e1530570aee5de086b3/Logic/PrimitiveRecursiveArithmetic/ConciseIPRA/}
\setlength{\emergencystretch}{3em}
\setcounter{tocdepth}{1}
\setlist{nosep}
\pagestyle{fancy}\fancyhf{}
\fancyhead[L]{\small Concise iPRA: checkers and progress}
\fancyhead[R]{\small 12 November 2026}
\fancyfoot[C]{\thepage}
\lstset{basicstyle=\ttfamily\scriptsize,columns=fullflexible,keepspaces=true,breaklines=true,showstringspaces=false,frame=single,rulecolor=\color{gray!35},aboveskip=0.8em,belowskip=0.8em}
\title{\textbf{iPRA: smaller checkers and\\the independence of progress}\\[0.6em]\large Two measured improvements and a coherent\\nonstandard model of the remaining kernel}
\author{Prepared for Vladimir Reshetnikov}
\date{12 November 2026}
\begin{document}
\maketitle
\begin{abstract}
Two complete compact checking programs improve the current repository's empirical size envelope: one uses 251 significant Python tokens instead of 253; another uses 502 normalized bytes instead of 506. They are different programs, and neither improves the incumbent 275-node AST record. These are implementation refinements of the existing bi-interpretable theory, not a new smaller arithmetic axiom basis.

The principal mathematical result answers the repository's open progress-independence question. A staged tower of ultrapowers supports every finite scalar iterator body, with actual equality and coherent substitution, while $I(c,0;u.u)=1\ne0$ at a nonstandard count. All seven non-progress kernel rules are sound in the resulting model. Nevertheless the literal equation $E(y):=I(y,0;u.u)=0$ has a provable base and successor step, so its progress closure is not derivable without the eighth rule. The construction can additionally preserve the ordinary commutative-monoid laws for addition and $I(n,0;u.Su)=n$. A separate finite-clock argument excludes every fixed finite successor cycle without induction.

The archive includes explicit rules, both ordinary interpretation composites, complete checker source, frozen comparisons, provenance, finite regression evidence, and reproducible LaTeX. No global minimality or new proof-assistant verification is claimed; the requested single construction below every existing size coordinate is not established.
\end{abstract}
\tableofcontents
\newpage

\section{Results and the exact comparison}
\label{sec:results}

The source snapshot is commit
\begin{center}\texttt{6ffc49f55d355b548f7a8e1530570aee5de086b3}.\end{center}
It contains ninety research packages and three unfiled archives in
\file{docs/incoming/}. The latter must be inspected: the top-level package
table alone misses some of the best implementations. The prior source
audit in \file{ipra_orbit_closure.zip} covers exactly the current Research
tree, whose Git tree hash is
\begin{center}\texttt{76de155aec9da910a74320fd692ed9ed790515c9}.\end{center}
The old audit commit and the present commit have this identical tree;
the new intake archives were separately scanned~\cite{repo,incoming}.

All measurements below use CPython 3.12.14, whole standalone modules,
\texttt{ast.unparse} normalization and one final newline. Significant tokens
exclude comments and layout; AST counts include every node of the reparsed
normalized source. The measured checkers contain no docstrings. Imports,
input handling, verification, and the exact success output \texttt{ACCEPT}
remain inside the measurement.

\begin{center}\small
\begin{tabular}{lrrr}
\toprule
Complete program & Tokens & AST nodes & Normalized bytes\\
\midrule
Frozen incoming token control & 253 & 280 & 599\\
\textbf{New token profile} & \textbf{251} & 283 & 605\\
Frozen repository byte control & 267 & 303 & 506\\
New parity companion & 270 & 303 & 503\\
\textbf{New byte profile} & 269 & 301 & \textbf{502}\\
New strict public-input companion & 421 & 471 & 869\\
\bottomrule
\end{tabular}
\end{center}

The old coordinate envelope is $(253,275,506)$; the new envelope is
$(251,275,502)$. These tuples combine different programs. In particular,
the 275-node program remains an inherited implementation: no supplied new
program has all three new scores. The byte improvement is $4/506$, about
$0.79\%$, and the token improvement is $2/253$, also about $0.79\%$.
These modest changes must not be described as a mathematical axiom deletion.

The observed envelope is an attained upper bound on the unknown minimum
implementation size. It is not a theorem excluding smaller programs or
different proof languages. The request to beat a ``lower bound'' is therefore
interpreted as a request to beat this empirical benchmark. We improve two
coordinates separately; we do not establish a universally simplest
axiomatization or a program below all three coordinates.

\paragraph{The mathematical advance.}
The model in Section~\ref{sec:progress-independence-audit} proves that
progress is independent of the other seven native rules. Together with the
seven existing independence results, this completes the irredundance
analysis of the original eight-rule kernel~\cite{models}. Irredundance
concerns literal deletion of a rule from that kernel. It does not preclude
the already successful rule fusions, changes of language, or alternate
bi-interpretations. Primitive recursion equations without induction have a
substantial earlier literature; Friedman studied their equational
consequences and the effect of successor coverage~\cite{friedman}. The
specific claim here is the proved countermodel for the repository's full
native scalar syntax, including its substitution requirements; no
literature-wide priority claim is made.

\section{A complete proof presentation and its arithmetic meaning}
\label{sec:presentation}

\subsection{Syntax and all eight primitive rules}

There is one numerical sort. The terms are
\[
t::=x\mid0\mid S(t)\mid I(n,a;u.h).
\]
Only $u$ in $h$ is bound by $I$; neither $n$ nor $a$ is within that binding.
Every body is finite earlier syntax. Equivalently, one can use a staged
first-order signature with an operation $I_h(n,a,\vec p)$ for each fixed
body $h(u,\vec p)$. This does not introduce program-valued variables or a
uniform total evaluator for all primitive recursive descriptions.

Native formulas and their ordinary readings are
\[
A::=(s=t)\mid Q_y A\mid D_{x,y}A,
\qquad
Q_yB\equiv B(0)\to\forall yB(y),
\]
\[
D_{x,y}B\equiv\exists x\bigl(B(x,0)\to\forall yB(x,y)\bigr).
\]
Contradiction is designated by $0=S0$. Contexts are finite lists of retained
native formulas. Capture-avoiding substitution and fresh variables are
understood. The precise eight-rule kernel $\K$ is the one of
\file{Lean/ConciseIPRA/Native.lean}~\cite{native}:

\begin{longtable}{p{0.13\textwidth}p{0.80\textwidth}}
\toprule
Rule & Entire schema\\
\midrule
\endhead
Hyp & $A\in\Gamma$ implies $\Gamma\vdash A$.\\[5pt]
C0 & $\Gamma\vdash I(0,a;u.h)=a$.\\[5pt]
CS & $\Gamma\vdash I(Sn,a;u.h)=h[I(n,a;u.h)/u]$.\\[5pt]
Tr & From $\Gamma\vdash s=t$ and $\Gamma\vdash h(s)=r$, infer $\Gamma\vdash r=h(t)$, for any one-term template $h$.\\[5pt]
QE & From $\Gamma\vdash Q_yB(y)$ and $\Gamma\vdash B(0)$, infer $\Gamma\vdash B(t)$.\\[5pt]
DE & From $\Gamma\vdash D_{x,y}B$ and $\Gamma,Q_yB(b,y)\vdash A$, infer $\Gamma\vdash A$; $b$ is fresh for the retained data and conclusion.\\[5pt]
DI & From $\Gamma,B(t,0)\vdash B(t,Sz)$, infer $\Gamma\vdash D_{x,y}B$; $z$ is fresh for the context, witness $t$, and passive parameters.\\[5pt]
QI & From $\Gamma,E(z)\vdash E(Sz)$, infer $\Gamma\vdash Q_yE(y)$; $E$ is a literal equation and $z$ is fresh for the retained context and passive parameters.\\
\bottomrule
\end{longtable}

Let $\Km$ be these same rules with QI deleted. No primitive ordinary
logical, successor-injectivity, reflexivity, or further induction axiom is
silently added. Weakening, exchange, contraction, cut, and the needed
equational rules are admissible. For example, let $d=I(0,t;u.u)$.
C0 gives $d=t$; Tr with the identity template and two copies of this
equation gives $t=t$. Tr with a constant template then gives symmetry,
and the identity template gives transitivity after one symmetry. These
constructions use no progress rule.

\subsection{The ordinary logical interface is explicit}

For clarity, the ordinary connectives need not be taken on trust as
additional rules. Write $C(v,a,b)=I(v,a;u.b)$, so that $C(0,a,b)=a$ and
$C(Sv,a,b)=b$. After replacing every auxiliary $Q$ by a $D$ with unused
witness, define a formula gate by
\[
G_v(s=t):=(C(v,0,s)=C(v,0,t)),\qquad
G_v(D_{x,y}B):=D_{x,y}G_vB.
\]
Renaming prevents capture. The inherited case-rule derivations give
$G_0A$, $G_{Sv}A\dashv\vdash A$, and
$A\dashv\vdash D_{r,y}G_yA$ when $r,y$ are fresh.

A selector between equations is
$M_v(s=t,r=w):=(C(v,s,r)=C(v,t,w))$.
For compound inputs, align their fresh outer binders and define
\[
\begin{aligned}
M_v(D_{x,y}A,D_{x,y}B)&=D_{x,y}M_v(A,B),\\
M_v(E,D_{x,y}B)&=D_{x,y}M_v(G_yE,B),\\
M_v(D_{x,y}A,E)&=D_{x,y}M_v(A,G_yE),
\end{aligned}
\]
where $E$ is an equation. This recursion decreases the maximum binder
depth and has $M_0(A,B)\dashv\vdash A$ and
$M_{Sv}(A,B)\dashv\vdash B$.
Put $\delta(v)=C(v,1,0)$ and $p(z)=I(z,0;u.\delta(u))$.
The complete macro inventory is
\[
\begin{aligned}
A\Rightarrow_*B&:=D_{r,z}M_z(A,B),\\
A\wedge_*B&:=D_{r,z}G_zM_{p(z)}(A,B),\\
\exists_*x\,A(x)&:=D_{x,y}G_yA(x),\\
A\vee_*B&:=\exists_*x\,M_x(A,B),\\
\forall_*x\,A(x)&:=A(0)\wedge_*D_{r,z}G_zA(z),\\
\bot_*&:=(0=1).
\end{aligned}
\]
Here the newly bound variables are fresh and $r$ is unused. The exact
derivations, including both formula round trips, are preserved in
\file{reference/foundations/inherited_choice_basis.tex}. This interface
uses cases and computations, rather than progress for quantified formulas.
QI remains restricted to literal equations; making it available for all
matrices would change the intended arithmetic strength.

\subsection{Both interpretations and both composites}
\label{sec:bridge}

The ordinary realization $T_I$ is intuitionistic first-order logic with
actual equality over the scalar terms. Its nonlogical axioms are C0, CS,
\[
\forall x\,(x=0\lor\exists y\ x=Sy),\qquad 0=S0\to\bot,
\]
and contextual induction for literal equations. The genuine falsity bridge
is explicit: positive computation equations alone have a singleton model.
The native kernel is a proof presentation of this ordinary theory. The
readings of $Q,D$, the logical macros above, and equation induction obtained
from QI followed by QE supply translations in both directions~\cite{bridge}.

The established arithmetic bridge has the following noncircular order.
First construct a scalar predecessor $P$ and prove $P0=0$ and $P(Sn)=n$.
Next construct a term reflecting actual equality into a Boolean zero test,
and compile quantifier-free formulas into characteristic equations. Only
then use the resulting quantifier-free induction to prove pairing and
projection identities, and simulate general primitive recursion. The
unchanged predecessor, reflection, and compilation proofs accompany the
archive; they are inherited results, not new claims of this report.

For example, if $\pi,L,H$ are the proved pairing and projections and
\[
f(0,\vec p)=g(\vec p),\qquad
f(Sn,\vec p)=h(n,f(n,\vec p),\vec p),
\]
the scalar compiler sets
\[
V(n,\vec p)=I\bigl(n,\pi(0,g^*(\vec p));
u.\pi(SL(u),h^*(L(u),H(u),\vec p))\bigr),
\quad f^*(n,\vec p)=H(V(n,\vec p)).
\]
Equation induction proves $L(V(n,\vec p))=n$; the projection identities
then prove both defining equations of $f^*$. Each finite description is
compiled externally; there is no single internal evaluator.

\begin{theorem}[Inherited bi-interpretation, preserved by the refinements]
The ordinary realization checked by the new programs is bi-interpretable
and equiconsistent with full first-order $\ipra$. Both numerical domains
are whole, equality is actual, and the graph of each numerical composite
isomorphism is $x=y$.
\end{theorem}
\begin{proof}
Interpret each scalar iterator in $\ipra$ by primitive recursion whose
update ignores its counter. This preserves its equations, coverage,
separation, and equation induction. In the other direction use the scalar
compiler and the predecessor and equality-reflection constructions just
specified; they recover the successor axioms and every quantifier-free
induction instance, including arbitrary retained contexts.

To check the composites, induct externally on finite operation
descriptions. After lower-level operations have been identified, an
original operation and its double interpretation satisfy the same initial
value and recurrence. Equation induction proves equality at every count.
This proves internal uniqueness and identifies their values over the
entire domain. Formula composites follow by structural induction from
the stated readings and the logical-macro round trips. The identity graph
is thus total, bijective, and preserves all interpreted operations.
The explicit falsity bridge transfers proofs of contradiction in both
directions.

The implementation refinements below preserve exactly the same canonical
public theorem set, by exact trace equivalence or canonical proof
compilation. Composing those transformations with the inherited
interpretations proves the assertion for the new checkers. No theory is
claimed to prove its own consistency.
\end{proof}


% Standalone inclusion fragment; requires amsmath, amssymb, amsthm.
% Independent audit of the staged-ultrapower construction.
\section{Independence of progress: a coherent model of scalar computation}
\label{sec:progress-independence-audit}

\begin{theorem}[Progress is not admissible]
The kernel $\Km$ does not derive
\[
 Q_y\bigl(I(y,0;u.u)=0\bigr).
\]
Its base instance and successor-step sequent are derivable in $\Km$.
There is a classical countermodel with actual equality, successor
injectivity, coverage, separation, and every finite scalar computation.
Consequently QI is independent of the other seven native rules.
\end{theorem}

The following proof takes place in classical metamathematics with choice.
Its conclusion is a countermodel to an intuitionistic calculus; using a
classical structure as an intuitionistic countermodel is sufficient.
No claim of a smaller complete axiomatization follows from this result.

\subsection{The semantic issue that must be addressed}

An arbitrary choice of an orbit for each \emph{written} update body is not
enough. Substituting terms for the passive parameters can change the written
body without changing its unary function. We therefore build an
\emph{extensional family} of scalar iterators. This is a stronger semantic
requirement than merely imposing each displayed computation equation, and
it prevents a gap in the native substitution lemma.

A body descriptor is a finite first-order term $h(u,\vec p)$ with one
distinguished state variable $u$ and a finite tuple of passive variables.
When an iterator for this descriptor is present, write its function symbol
as $I_h(n,a,\vec p)$. Its computation axioms are
\begin{align}
 I_h(0,a,\vec p)&=a,\label{eq:audit-c0}\\
 I_h(Sn,a,\vec p)&=h(I_h(n,a,\vec p),\vec p).
 \label{eq:audit-cs}
\end{align}
For \emph{every pair} of available descriptors, including two copies of the
same descriptor with different parameter values, require the first-order
schema
\begin{equation}
\begin{split}
 \forall\vec p\,\forall\vec q\bigl[
  (\forall u\ h(u,\vec p)=g(u,\vec q))\ \longrightarrow {}&\\[-0.6em]
  \forall n\,\forall a\ I_h(n,a,\vec p)=I_g(n,a,\vec q)
 \bigr].&
\end{split}
\tag{Ext}\label{eq:audit-ext}
\end{equation}
Thus an available iterator depends on the actual unary function and the
two numerical inputs, rather than on a particular presentation of that
function. Parameter repetition, omission, and reordering are included.

\subsection{The orbit supplied by one ultrapower}

Let $A$ be any nonempty structure, let $U$ be a nonprincipal ultrafilter on
$\mathbb N$, and put $B=A^{\mathbb N}/U$. Fix a unary term function of $B$,
\[
 f(v)=h^B(v,\vec p),\qquad \vec p=[\vec p_i],
\]
and an element $a=[a_i]$ of $B$. Write
$f_i(v)=h^A(v,\vec p_i)$. For each ordinary integer $j\in\mathbb Z$, set
\[
 x_j=\bigl[f_i^{\,\max(0,i+j)}(a_i)\bigr].
\]
For all sufficiently large $i$,
$\max(0,i+j+1)=\max(0,i+j)+1$. Hence
\[
                         f(x_j)=x_{j+1}
                         \quad(j\in\mathbb Z).
\]
Every parameter-definable unary term function in this ultrapower therefore
has a full two-sided orbit. The iterations in the coordinates are finite
external iterations; no arithmetic induction axiom in $A$ is being used.
The displayed sequence need not pass through $a$, and no such condition is
required on a successor component that does not contain zero.

\subsection{Extension of a coherent partial family}

Suppose $A$ has successor injectivity, zero outside the image of successor,
zero/successor coverage, and no finite successor cycles. Suppose it also has
a collection of iterator symbols satisfying
\eqref{eq:audit-c0}, \eqref{eq:audit-cs}, and all instances of
\eqref{eq:audit-ext} for that collection.

Take the ultrapower $B$ in the language containing \emph{all} those symbols.
All the stated first-order conditions transfer by \L o\'s's theorem.
In particular, the old iterator values are retained, and their coherence
holds also for parameter tuples newly present in the ultrapower.

The successor graph of $B$ consists of the standard ray
$0,S0,S^2 0,\ldots$ and components isomorphic to $\mathbb Z$.
Indeed, every nonzero point has a unique predecessor. A point outside the
standard ray cannot reach zero after finitely many predecessor steps;
injectivity and the absence of finite cycles ensure that all steps in its
two-sided component are distinct. Choose an origin in each such component.

Now add iterators for any set of descriptors in the \emph{old language}.
Partition their parameter instances by the actual unary functions
$f:B\to B$ that they define.

If a function $f$ already has an iterator from the old family, use that
iterator at every counter and every seed. This prescription is well-defined:
if $f=h^B(-,\vec p)=g^B(-,\vec q)$ for two old descriptors, Ext says that
their iterator values agree for all counters and seeds.

For a function not represented in the old family, make one simultaneous
choice for its entire equivalence class. At a standard counter $S^k0$ put
\[
                         I_f(S^k0,a)=f^k(a).
\]
On a nonstandard successor component with chosen origin $b$, choose a
two-sided orbit $(x_j)_{j\in\mathbb Z}$ supplied by the preceding lemma,
and set
\[
                         I_f(S^j b,a)=x_j.
\]
Here negative powers mean the unique iterated predecessor on that component.
Choices may be made separately for each $f$, seed $a$, and component.
The set-theoretic axiom of choice suffices for all these choices.

Define every newly added $I_h(n,a,\vec p)$ by the chosen iterator for the
actual function $h^B(-,\vec p)$. Its two computation equations hold by
construction. Ext holds for every pair of old or new symbols, because
equal unary functions were assigned the same entire iterator. The expansion
changes no old function or old element.

\subsection{A seed model in which identity iteration is not constant}

Start with the disjoint union
\[
                         A_0=\mathbb N\ \sqcup\ \mathbb Z.
\]
Zero belongs to the first summand; successor is ordinary successor on each
summand. Let $1=S0$. Interpret the iterator $J=I_u$ of the identity update by
\[
 J(n,a)=
 \begin{cases}
 a,&n\in\mathbb N,\\
 1,&n\in\mathbb Z.
 \end{cases}
\]
Then $J(0,a)=a$ and $J(Sn,a)=J(n,a)$. All the required successor conditions
hold. The initial extensionality condition for the sole descriptor $u$ is
automatic. For any chosen $c$ in the $\mathbb Z$ summand,
\[
                           J(c,0)=1\ne0.
\]
No predicate distinguishing the two summands is required in the language.

Iterate the extension construction through $\omega$ stages. At each stage,
take an ultrapower in the entire current language and add iterators for all
finite unary term descriptors of that language. Identify the preceding
structure with its diagonal elementary image before adding the symbols.
The languages and structures can then be regarded as increasing chains.
The reduct at every later stage is an elementary extension in each earlier
language. Their union $A_\omega$ retains all computation equations, all
cross-extensionality schemas, all successor conditions, and the distinguished
equation $J(c,0)=1\ne0$.

To justify the union step explicitly, each fixed old-language sentence is
handled by the elementary-chain theorem on the reducts to that language.
Every computation or Ext instance uses only finitely many symbols and was
imposed at a finite stage, so it holds in the union. Equivalently, Ext can
be checked directly: if two final unary functions agree, they agree on
every sufficiently late stage containing their parameters; applying Ext at
a stage also containing the proposed counter and seed gives equality of
the two iterator values there and hence in the union.

\subsection{Returning to native scalar syntax}

Every finite native term
\[
 t::=x\mid0\mid St\mid I(n,a;u.h)
\]
has an interpretation using finitely many staged symbols. Expand its inner
iterators first. The translated update body is then a finite term of some
stage language, so its descriptor has an iterator at the next stage.
Parameters are ordinary numerical variables; there is no function sort and
no internal universal evaluator.

More intrinsically, for any unary function definable by a finite native
body with parameters in $A_\omega$, choose a staged descriptor presenting
it and use its iterator. Ext makes this choice independent of the
descriptor. It follows by structural induction on native terms that
capture-avoiding substitution has the usual semantics:
\[
 \llbracket t[s/x]\rrbracket_\rho
 =\llbracket t\rrbracket_{\rho[x\mapsto\llbracket s\rrbracket_\rho]}.
\]
At an iterator, the induction hypothesis identifies its counter, seed, and
pointwise update function; Ext identifies the resulting iterators. This is
the exact point at which mere independent choices for written bodies would
have been insufficient.

\subsection{Soundness without progress and the refuted instance}

Interpret native equality as actual equality, and use the ordinary readings
\[
 Q_y B=B(0)\to\forall y\,B(y),\qquad
 D_{x,y}B=\exists x\bigl(B(x,0)\to\forall y\,B(x,y)\bigr).
\]
Hypothesis, equality transport, the two computations, closure elimination,
and witness elimination are sound. Case introduction is sound because every
element is zero or a successor: from $B(t,0)$ and all successor instances
$B(t,Sz)$ one obtains every instance. The usual freshness conditions ensure
that $t$ and the retained assumptions do not depend on the fresh $z$.
All these arguments are ordinary first-order soundness arguments in the
single structure $A_\omega$.

Let
\[
                          E(y)\;:\equiv\;I(y,0;u.u)=0.
\]
The zero computation proves $E(0)$. The successor computation and ordinary
derived equality reasoning prove
\[
                          E(z)\ \vdash\ E(Sz)
\]
without any induction, since
$I(Sz,0;u.u)=I(z,0;u.u)=0$ under the displayed hypothesis.
Nevertheless $E(c)$ is false in $A_\omega$. Consequently $Q_yE(y)$ is
false there and is not derivable by the kernel without progress.
Thus QI by progress is not admissible over the other seven rules.
With an unused existential witness, the same model also refutes the
witnessed progress conclusion $D_{x,y}E(y)$.

\paragraph{Proof status and limits.}
This is a written model construction, not a proof-assistant verification.
The essential requirements are simultaneous extensional choice for actual
unary functions; universal cross-extensionality for all parameter tuples;
ultrapowers in the full current language; and adding only prior-language
body descriptors at a given stage. Omitting any of those requirements would
leave a genuine coherence or circularity gap. The result proves independence
of progress. It does not prove a complexity lower bound for other
presentations, and it does not satisfy a requirement for a smaller complete
proof checker by itself.

% Optional strengthening of the ultrapower independence construction.
\subsection{A stronger initial structure with cancellative addition}

The initial structure can already satisfy additive cancellation and the
identity law for successor iteration. Take
\[
 M_0=\{(0,n):n\in\mathbb N\}
     \cup\{(k,z):k\in\mathbb N_{>0},\ z\in\mathbb Z\},
\]
with componentwise addition, $0=(0,0)$, $1=(0,1)$, and
$S(k,z)=(k,z+1)$. It is a submonoid of the additive group $\mathbb Z^2$,
so addition is associative, commutative, and cancellative, with identity
$0$. It is conical: $a+b=0$ implies $a=b=0$. Successor is injective,
omits zero, has no finite cycle, and every nonzero element has a
predecessor in the domain.

Preinstall two iterator symbols:
\[
 J(n,a)=
 \begin{cases}
 a&\text{if the first coordinate of }n\text{ is }0,\\
 1&\text{otherwise},
 \end{cases}
 \qquad K(n,a)=a+n.
\]
They represent the updates $u\mapsto u$ and $u\mapsto Su$, respectively.
Both satisfy their zero and successor computations. The added identity
law $K(n,0)=n$ holds, as do the usual additive zero and successor laws.
Their cross-extensionality condition has a false antecedent, since
identity and successor are distinct functions. Extensionality within
each description is immediate. Thus this is a valid starting point for
the same extension tower. All these first-order laws survive that tower.

At $c=(1,0)$, however, $J(c,0)=1$. Hence the final countermodel can have
all of the listed additive laws, including cancellation, and the
successor identity orbit, while still refuting literal zero-equation
induction.

This does not furnish a countermodel to full successor uniqueness.
Indeed, already in the initial structure the term
\[
 f(n)=n+J(n,0)
\]
satisfies $f(0)=0$ and $f(Sn)=Sf(n)$, whereas
$f(c)=Sc\ne c=K(c,0)$. Its addition is itself the successor-iterator
operation $K$, so this failure uses scalar terms. Full successor
uniqueness therefore fails in this structure, as it must.

% A new mathematical supplement. No claim of checker-size improvement.
% Requires amsmath, amssymb, amsthm; use own theorem environments if needed.
\section{Finite scalar clocks without induction}
\label{sec:finite-clocks}

This section resolves the even-cycle gap in the linked README's first
further question as a written theorem. It does not remove the progress rule or establish a smaller
complete checker. The result is a constraint on any proposed countermodel
for the independence of progress.

Work first in intuitionistic first-order logic with actual equality, the
two scalar iterator computations, and numerical coverage
\[
 x=0\ \lor\ \exists y\,x=Sy.
\]
There is no induction rule. Write
\[
 C(x,a,b)=I(x,a;u.b),\qquad Z(x)=C(x,1,0),\qquad
 \bar j=S^j0.
\]
Thus $C(0,a,b)=a$, $C(Sx,a,b)=b$, and coverage proves
$Z(x)=0\lor Z(x)=1$. All inductions below are external inductions on a
fixed finite integer specifying a term, not induction in the object theory.
For a fixed unary term $h$, $h^{[r]}$ means $r$-fold syntactic composition;
it does not introduce a new iterator or an object-language quantifier.

\begin{lemma}[Finite clocks]
For every fixed integer $m\geq2$, there is a finite scalar term $h_m(x)$
such that the theory without induction proves
\[
 h_m(x)\in\{\bar0,\ldots,\overline{m-1}\},\qquad
 h_m(\bar j)=\overline{j+1}\quad(0\leq j<m-1),\qquad
 h_m(\overline{m-1})=0.
\]
Here membership abbreviates a finite disjunction of actual equalities.
\end{lemma}
\begin{proof}
Start with $h_2=Z$. Given $h_m$, define
\[
 T_m(x)=I(x,0;u.h_m(u)).
\]
Coverage on $x$ and the two computations show that $T_m(x)$ belongs to
$\{\bar0,\ldots,\overline{m-1}\}$: the zero case is $0$, and in the
successor case its value is an $h_m$-value. No assertion about the preceding
iterator value is needed.

For $1\leq r<m$, put
\[
 d_{m,r}(z)=Z\bigl(h_m^{[m-r]}(z)\bigr),\qquad
 V_{m,0}(z)=0,
\]
\[
 V_{m,r}(z)=C\bigl(d_{m,r}(z),V_{m,r-1}(z),\overline{r+1}\bigr),
 \qquad V_m=V_{m,m-1}.
\]
Finite ground computation proves
$d_{m,r}(\bar s)=1$ for $s=r$, and $d_{m,r}(\bar s)=0$ for
$0\leq s<m$, $s\ne r$. Consequently
\[
 V_m(0)=0,\qquad V_m(\bar r)=\overline{r+1}\quad(1\leq r<m).
\]
Every value of $V_m$, at an arbitrary argument, is one of
$0,\bar2,\ldots,\bar m$, by finitely many applications of coverage to
the selector counters. Define
\[
 \boxed{\quad h_{m+1}(x)=C\bigl(x,1,V_m(T_m(x))\bigr).\quad}
\]
Its range is $\{0,\bar1,\ldots,\bar m\}$. At $x=0$ its value is $1$.
At $x=\bar j$ with $1\leq j<m$, finite computation gives
$T_m(\bar j)=\bar j$, so its value is $\overline{j+1}$. At $x=\bar m$,
$T_m(\bar m)=0$, so its value is $0$. This completes the external
construction.
\end{proof}

The first nontrivial extension has the compact description
\[
 h_3(x)=C\bigl(x,1,C(Z(Z(T_2(x))),0,2)\bigr),
 \qquad T_2(x)=I(x,0;u.Z(u)).
\]
Direct computation gives $h_3(0)=1$, $h_3(1)=2$, and $h_3(2)=0$.

\begin{lemma}[Separation of finitely many numerals]
For fixed $m\geq2$ and distinct $i,j<m$, the assumption
$\bar i=\bar j$ proves $0=1$, without induction or successor injectivity.
\end{lemma}
\begin{proof}
Apply $h_m^{[r]}$ with $r\equiv-i\pmod m$ and $0\leq r<m$.
The resulting equation is $0=\bar s$ for some $1\leq s<m$.
Applying $Z$ gives $1=0$, whose symmetric form is the desired
contradiction. Only finitely many computation equations are used.
\end{proof}

\begin{theorem}[No finite successor cycle without induction]
For every fixed integer $k\geq1$ and every scalar term $c$, the theory
without induction derives
\[
 S^k c=c\ \longrightarrow\ 0=1.
\]
With $0\ne1$ it therefore proves $S^k c\ne c$.
\end{theorem}
\begin{proof}
Choose $m=k+1$ and set $v=T_m(c)$. The hypothesis and $k$ applications of
successor computation give
\[
 v=T_m(c)=T_m(S^k c)=h_m^{[k]}(v).
\]
The finite-range property of $T_m$ gives the cases $v=\bar j$ for
$0\leq j<m$. In each case the displayed equation becomes
\[
 \bar j=\overline{j+k\pmod m}.
\]
Since $k=m-1$, the two residues differ. The preceding separation lemma
gives $0=1$ in every case.
\end{proof}

\subsection{Removing the finite disjunctions in the native calculus}

There is also a direct written translation into the native kernel without
progress, so the preceding use of ordinary disjunction is dispensable.
The existing theorem \texttt{NoInd.cases} derives the following rule for
equation templates: if $E(0)$ and $E(Sz)$ are derivable uniformly in the
retained context, then $E(t)$ is derivable for every term $t$. The
computation and crossed equality rules provide ordinary equational
congruence and transitivity without progress.

Call a term a finite selector tree over $L$ if it is built from leaves in
the fixed finite list $L$ by $C(q,a,b)$, with arbitrary counter terms $q$.
For any equation template $E$, if $E(\ell)$ is derivable uniformly for
each leaf $\ell\in L$, then $E(t)$ is derivable for every such tree $t$.
Prove this externally by induction on its finite syntax: at a selector,
apply \texttt{NoInd.cases} to its counter and compute the two branches.
Each $h_m(x)$ is literally a finite selector tree over
$L_m=\{\bar0,\ldots,\overline{m-1}\}$. A further application of
\texttt{NoInd.cases} to the counter of $T_m(c)$ shows:
\[
 \bigl(E(\bar j)\text{ derivable for every }j<m\bigr)
 \quad\Longrightarrow\quad E(T_m(c))\text{ derivable}.
 \tag{FR}
\]
The zero branch is $E(0)$; the successor branch is
$E(h_m(T_m(z)))$, handled by the selector-tree argument.

Put $b_r(a)=Z(h_m^{[(-r\bmod m)]}(a))$ for $0\leq r<m$.
These are Boolean terms, with $b_r(\bar j)=1$ exactly when $r=j$.
For Boolean arguments define
\[
 \operatorname{xor}(p,q)=C(p,q,Z(q)),\qquad
 \operatorname{or}(p,q)=C(p,q,1).
\]
Let $D_m(a,b)$ be the left-associated finite disjunction (starting at
$0$) of the terms $\operatorname{xor}(b_r(a),b_r(b))$.
Finite zero/successor cases prove $D_m(a,a)=0$ for every $a$; no
disjunction formula or equality characteristic for arbitrary numbers is
used. For $k=m-1$, ground computation proves
\[
 D_m(\bar j,h_m^{[k]}(\bar j))=1\qquad(j<m).
\]
Apply (FR) to this equation template to obtain
$D_m(v,h_m^{[k]}(v))=1$ at $v=T_m(c)$. The cycle assumption gives
$v=h_m^{[k]}(v)$, so equational transport and $D_m(v,v)=0$ give
$0=1$. This is the required native derivation, with no progress rule.

\paragraph{Boundary of the result.}
The construction and the native translation are written proofs; no new
Lean verification is claimed. They do not prove the progress rule
redundant, and do not by themselves settle whether successor uniqueness
is equivalent to full iPRA. The novelty is the general finite-clock
construction, which completes the existing odd-cycle exclusion.

The companion \texttt{check\_finite\_clocks.py} builds genuine scalar
terms using de Bruijn variables and evaluates them on bounded standard
inputs. It stores repeated subexpressions as a directed acyclic graph;
this is only a representation of a finite expanded term. It represents
$h_m^{[r]}(a)$ as $I(\bar r,a;u.h_m(u))$, which expands to the same
composition by finitely many computations. Its default run checks
periods $2$ through $12$ and numerical inputs $0$ through $32$, including
finite range, clock values, cycle transitions, and the Boolean
discriminator: 1,243 checks pass. These evaluations corroborate finite
examples only; the external construction and proofs above establish the
universal theorem.

\section{Smaller implementations of the complete calculus}
\label{sec:checkers}

\subsection{The measured public contract}

The compact programs take a finite JSON certificate consisting of integers
and arrays, with a canonical closed guarded Eq/$D$ public goal. Intermediate
proofs use the inherited larger auxiliary language. A guard replaces an
equation $s=t$ below a binder by
\[
I(y,0;u.s)=I(y,0;u.t),
\]
under a closure with unused witness. Its zero instance is reflexive; its
successor instances express $s=t$. This is why restricting public atomic
matrices to guarded ones preserves the ordinary theory while permitting a
smaller root-selected introduction~\cite{public}.

The compact modules do not themselves reject every malformed public
grammar or every non-integer JSON leaf. The 869-byte strict companion
includes those checks, including unused certificate prefixes. It is a
separate program under a stronger input contract. A strict program's size
must not be compared as though it had the same job as the compact
506-byte control. All programs import only the Python JSON library;
no arithmetic oracle, local checker module, embedded compressed program,
\texttt{eval}, or \texttt{exec} is used.

The inherited arity codec stores a term or formula as an integer/list
tree, using weighted binding slots. Its common opener is denoted below by
$J_r(t)$: opening tree $t$ with image $r$ at formula depth zero and the
current numerical scope. Its recursive depth increment is $5-\ell$, where
$\ell$ is the node arity. Canonical numerical constructors have odd arity
$1$ or $3$. Equations have arity $2$, paired formulas arity $1$, and auxiliary
closures arity $3$. Full decoding and compilation specifications are included
in \file{reference/foundations/inherited_arity_checker.tex} and
\file{reference/COMPACT_CONTRACT.txt}.

\subsection{The 251-token refinement is trace-equivalent}

The 253-token control reads a four-field instruction as
\begin{lstlisting}[language=Python]
m, s, u, t = q()
\end{lstlisting}
The new program reads
\begin{lstlisting}[language=Python]
m, *r = q()
s, u, t = r
\end{lstlisting}
and reuses $r$ where the control constructs \texttt{[s,u,t]}, and
\texttt{r[:2]} where it constructs \texttt{[s,u]}.
The second unpack enforces exactly three tail fields, so the required
instruction arity remains four. If both unpacks succeed, the variables
have exactly their former values and $r$ is exactly the former constructed
list. No scope, context, dispatch test, stream consumption, or inference
changes. If either unpack fails, both programs reject. Thus successful
traces agree for every finite JSON certificate, apart from finite runtime
resource limits. The measured saving is two significant tokens, with a
cost of three AST nodes and six normalized bytes relative to its matched
control. This is an implementation refinement, not a new inference rule.

\subsection{The 502-byte refinement and its completeness argument}

The 506-byte control already fuses continuation opening, allocation of a
new parameter, and context extension. Its continuation selector has the
form
\[
\mathcal C\mathbin{\texttt{ and }}J_s(t)
\mathbin{\texttt{ or }}a,
\quad
\mathcal C:=((m,a)=(M,J_u(t))),
\]
where $M$ is the required equality or closure major. The 503-byte parity
companion selects instead from the tuple $(a,J_s(t))$ with Boolean index
$\mathcal C$. This eagerly evaluates $J_s(t)$.

For integer/list data, whenever the eager evaluations succeed, the two
values coincide. If $\mathcal C$ is false both select $a$. If it is true
and $J_s(t)$ is truthy, both select $J_s(t)$. If $J_s(t)$ is false, then
the tree has an empty-list or scalar root: a successful opening of a
nonempty list is nonempty. At formula depth zero those empty/scalar
openings do not depend on their proposed image, so
$J_s(t)=J_u(t)=a$, and the selectors still agree. The eager companion
can reject additional noncanonical certificates, but cannot create a new
successful trace relative to the parity control.

Canonical completeness survives eager evaluation. Genuine elimination
already requires the source opening. In a computation packet the otherwise
inspected field is a numerical update body; in an introduction packet it
is a numerical witness. Starting at depth zero, their odd node arities
keep every encountered depth even. At an integer leaf the opener either
retains a legal coordinate or returns padding zero, without opening the
proposed image. Thus the extra evaluation terminates on every such
canonical field. Unused fields of witness-elimination packets can be set
to zero. This is a standardization of certificate data, not a restriction
of the theorems that have certificates.

The final byte replaces \texttt{\~{}len(a) \% 2} by
\texttt{3 \% len(a)}, as described in the incoming smaller-envelope
report. On the meaningful matrix arities the values are identical:
\begin{center}
\begin{tabular}{lccc}
\toprule
Matrix & $D$ (arity 1) & Equation (arity 2) & $Q$ (arity 3)\\
\midrule
Both mode values & 0 & 1 & 0\\
\bottomrule
\end{tabular}
\end{center}
An empty matrix now rejects by division by zero. Surplus arities in the
permissive administrative decoder mean the true equation $0=0$; a changed
mode there affects only a truth alias or the substitution instance of a
computation equation. No false public consequence is introduced. This last
argument uses the inherited sorted decoder, so the 502-byte program is
not claimed to have exactly the same raw acceptance language as the
control. The 503-byte parity companion is supplied for the cleaner trace
inclusion statement.

\paragraph{Freshness is retained.}
The expression \texttt{k := \~{}-k} is inherited from the 506-byte control
and means $k-1$. Source and target openings in continuation selection occur
before allocation; the outer opening and context extension see the new
scope. The new programs preserve this order. Replacing the scope counter
naively by context length would get that phase boundary wrong; none of the
delivered checkers makes that replacement.

\subsection{The strict companion}

After JSON parsing, the strict module first checks that the printed value
uses only spaces, brackets, commas, minus signs, and decimal digits.
On parsed JSON values, this character test accepts precisely integer/list
trees: strings introduce quotes; objects braces or colons; Boolean values
and null introduce letters; non-integral numbers introduce a decimal point
or letters. The public grammar validator then checks Eq/$D$ arities,
numerical constructors, the permitted numerical slot form, and literal
guards immediately below a paired binder. The kernel's initial opening
checks the numerical scope. These checks are inside the measured module,
not external preconditions of the strict companion.

\subsection{Executable evidence and its boundary}

The delivered \file{checker/reproduce.py} measures every whole module,
checks frozen source identities, and replays 512 checker processes under
ordinary and optimized Python. Every observed outcome matches its stated
expectation. Fixtures include genuine uses of an induction hypothesis,
arbitrary specialization, future-variable capture, context leaks,
false goals, surplus arities, erased redundant fields, and malformed
prefixes or packets. Global-scope token fixtures and local-scope byte
fixtures are separate because they use different certificate conventions.

The benchmark audit measures all 1,416 Python files in the three incoming
archives, including the 1,380 distinct frozen source blobs of the prior
Research audit. Programs below the previous byte record were inspected
as utilities, not standalone theorem checkers. The current Research tree
identity and all classification receipts are in
\file{checker/evidence/}. The record claim is limited to this inspected
corpus and its documented eligibility rules.

Finite tests verify concrete executions. The universal mathematical
arguments are the trace and compiler proofs above and the inherited
arithmetic bridge. The nonstandard countermodel is not simulated or
established by these tests. No new Lean or Rocq verification is asserted.

\section{What the new mathematics permits us to conclude}
\label{sec:conclusions}

The independence theorem closes a specific gap in the current research
program. The equation $I(y,0;u.u)=0$ is an explicit witness that all scalar
computations, equality, cases, and the ordinary $Q/D$ rules do not force
equation induction. The countermodel is acyclic, has actual equality, and
even admits a cancellative addition with the usual elementary identities.
It does not exploit an ill-typed formula, a failed successor cover, a
finite cycle, or incompatible interpretations of substituted bodies.

Combined with the seven independent deletions already established in
\file{NativeModels.lean}, this yields an eight-out-of-eight irredundance
result for the original rule inventory. The eight rules are not a global
minimum: the repository itself has equivalent presentations with fewer
rule families, obtained by fusion or a changed syntax. Likewise,
the smaller token and byte programs do not establish an absolute minimum
description length.

For the requested axiomatization, the honest outcome is therefore a
complete inherited bi-interpretable presentation with two smaller
implementations, together with a mathematical theorem showing why one
tempting deletion cannot work literally. A new mathematical axiom basis
with a single complete program below all three current coordinates
remains unconstructed.

\section{Further research questions}

\begin{enumerate}
\item Can one complete program beat 251 significant tokens, 275 AST nodes,
and 502 normalized bytes simultaneously under a fixed common input
contract? A joint result must be one measured program, not an envelope
assembled from several files.
\item Can the ultrapower extension lemma be formalized in the existing Lean
development using arbitrary carriers and interpreted successor? The current
natural-number readings fix the external standard successor and therefore
cannot witness the new independence result.
\item Can a different presentation combine computational equality with a
weaker induction interface without reintroducing its complexity in syntax
validation? The independence proof rules out deletion in the fixed kernel;
it leaves such a redesign open.
\item Which fragments of the coherent extension lemma survive if additional
iteration laws, such as fusion or uniqueness, are imposed? Computation and
extensionality of the update alone are compatible with the counterexample.
Uniqueness for the identity update immediately excludes its specific
witness, but excluding one model does not prove full arithmetic strength.
\item Can a counterpart of the new countermodel decide the remaining
successor-uniqueness versus full-equation-induction question? The strengthened
model below refutes full successor uniqueness and so is not a negative
answer to that different question.
\end{enumerate}

\appendix
\section{Transfer and the countable-model refinement}
\label{app:transfer}

For completeness, the model construction uses ordinary ultrapowers, not
set-theoretic large-cardinal ultrapowers. Fix a nonprincipal ultrafilter
$U$ on the ordinary natural numbers, available in the classical metatheory
with choice. Elements of $A^{\mathbb N}/U$ are sequences modulo agreement
on a set in $U$; operations are interpreted coordinatewise. For every
first-order formula $\varphi$,
\[
A^{\mathbb N}/U\models\varphi([\vec a_i])
\quad\Longleftrightarrow\quad
\{i:A\models\varphi(\vec a_i)\}\in U.
\]
Here is the needed proof outline. The atomic case is the definition of
the quotient and of coordinatewise operations. Finite Boolean connectives
use the ultrafilter laws. For an existential formula, a witness sequence
gives the forward direction; in the other direction choose a witness in
each coordinate where one exists, and an arbitrary element elsewhere.
This proves the statement by structural induction on formulas. In
particular the diagonal embedding by constant sequences is elementary.
This is the usual proof of \L o\'s's theorem.

The language at each stage is countable, and so is the union language.
Although the constructed union may be uncountable, it has a countable
elementary submodel containing the tracked element $c$. One may construct
such a submodel by starting with $c$ and the named constants, closing
under the countably many function symbols, then choosing witnesses for
every existential formula with parameters already selected, and repeating
this countably many times. The witness property gives an elementary
submodel by induction on formulas. It retains every computation and Ext
sentence as well as $J(c,0)=1\ne0$.

Thus the countermodel can be made countable. This does not contradict the
repository's obstruction to standard-natural-number readings: the
successor on this countable carrier has nonstandard components and is not
the externally fixed map $n\mapsto n+1$ after an arbitrary enumeration.
The proof does not claim a computable presentation of the full model.

\section{Complete compact checker listings}
\label{app:listings}

The following are complete programs, including input parsing and the
success signal. Their larger strict companion and frozen controls are
separate files in the archive. Acceptance means normal termination with
exact output \texttt{ACCEPT} and a newline; exceptions reject. As with
the comparison programs, the theoretical algorithms allow unbounded
finite resources, while a finite Python recursion or memory limit can
reject a sufficiently large certificate.

\subsection{251-token program}
\begin{lstlisting}[language=Python]
from json import load
k = 0

def i(a, t=0j, d=0):
    try:
        return [i(b, t, d + 5 - len(a)) for b in a]
    except:
        return a if k <= a < d else d % 2 and i(t, d=3)

def c(a, g):
    global k
    while i(a) not in g:
        m, *r = q()
        s, u, t = r
        if i(t, u) == a:
            m = r[:2] if s else [t] * 3
            a = i(t, s)
        h = c(m, g) or m * 3 if m else [[s and [s], u, t], i(t, r, 3) if s else u] if u else (a := i(*a, t, 2))
        k -= 1
        g += (i(h, (m or len(a) == 2) and k),)
        a = i(a, [k])
q = load(open(0)).pop
c(q(), ())
print('ACCEPT')
\end{lstlisting}

\subsection{502-byte program}
\begin{lstlisting}[language=Python]
import json

def c(a, g=[0], k=0):
    i = lambda a, t=c, d=0: [i(b, t, d + 5 - len(a)) for b in a] if a * 0 != 0 else a if k <= a < d else d % 2 and i(t, d=3)
    if i(a) in g:
        return
    m, s, u, t = q()
    h = c(m, g, k) or m * 3 or (u and [[s and [s], u, t], i(t, [s, u, t], 3) if s else u]) or (a := i(*a, t, 2))
    c(i((a, i(t, s))[(m, a) == (s and [s, u] or [t] * 3, i(t, u))], [(k := (~-k))]), g + [i(h, (m or 3 % len(a)) and k)], k)
q = json.load(open(0)).pop
c(q())
print('ACCEPT')
\end{lstlisting}

\begin{thebibliography}{9}
\bibitem{repo} V. Reshetnikov, \emph{ProveIt: Concise presentations of intuitionistic primitive recursive arithmetic}, commit \texttt{6ffc49f55d355b548f7a8e1530570aee5de086b3}, 12 November 2026. \href{\srcbase README.md}{Pinned project README}, especially Further questions 1.
\bibitem{native} ProveIt, \file{Lean/ConciseIPRA/Native.lean} and \file{NativeEq.lean}, same snapshot. \href{\srcbase Lean/ConciseIPRA/Native.lean}{Explicit eight-rule kernel}.
\bibitem{models} ProveIt, \file{Lean/ConciseIPRA/NativeModels.lean}, same snapshot. \href{\srcbase Lean/ConciseIPRA/NativeModels.lean}{Seven rule-independence models and the open progress case}.
\bibitem{bridge} ProveIt, \file{Research/certified_zero/reference/inherited_choice_basis.tex}, \file{inherited_ordinary_bridge.tex}, \file{scalar_predecessor.tex}, \file{equality_reflection.tex}, and \file{primitive_recursion.tex}, same snapshot. \href{\srcbase Research/certified_zero/reference/inherited_ordinary_bridge.tex}{Whole-domain arithmetic bridge}. Unchanged copies accompany this report.
\bibitem{public} ProveIt, \file{Research/deterministic_introduction/implementation/code/compact_optional.py}, Git blob \texttt{a0af7afb236b530eb30c63c5ab788151024271b7}, and the inherited guarded public-language proofs. \href{\srcbase Research/deterministic_introduction/implementation/code/compact_optional.py}{Frozen 506-byte comparator}.
\bibitem{incoming} ProveIt, \file{docs/incoming/ipra_smaller_envelope.zip}, \file{ipra_frontier_review_2026-11-12.zip}, and \file{ipra_orbit_closure.zip}, same snapshot. The first supplies the 253-token and 275-node coordinate records; the third audits the 506-byte source. \href{https://github.com/VladimirReshetnikov/ProveIt/tree/6ffc49f55d355b548f7a8e1530570aee5de086b3/Logic/PrimitiveRecursiveArithmetic/ConciseIPRA/docs/incoming}{Pinned intake directory}.
\bibitem{friedman} H. Friedman, \emph{Provable equality in primitive recursive arithmetic with and without induction}, Pacific Journal of Mathematics \textbf{57}(2), 379--392 (1975). \href{https://msp.org/pjm/1975/57-2/pjm-v57-n2-p10-s.pdf}{Original paper}.
\end{thebibliography}
\end{document}
